\documentclass[10pt,a4paper]{amsart}
\usepackage[margin=19mm]{geometry}
\usepackage{amsmath,amssymb,mathtools}
\usepackage[scr=rsfs,cal=euler]{mathalfa}
\usepackage{xfrac}
\usepackage[unicode,hidelinks,bookmarks=false]{hyperref}
\newcommand{\ie}{i.e.\ }
\newcommand{\cf}{cf.\ }
\newcommand{\resp}{respectively\ }
\newcommand{\Subsection}{\subsection}
\providecommand{\nobreakdash}{\nobreak\hbox{-}}
\setlength{\emergencystretch}{2em}
\multlinegap0pt
\usepackage{amssymb}
\usepackage{xcolor}
\newtheorem{lemma}{Lemma}
\title{A note on Galois groups of linearized polynomials}
\author{Peter M\"uller}
\date{Source: arXiv:2311.14953v1 (CC BY 4.0)}
\begin{document}
\maketitle

\noindent\emph{Source and licence.} A note on Galois groups of linearized polynomials. Version arXiv:2311.14953v1 (CC BY 4.0), distributed by arXiv under the Creative Commons Attribution 4.0 International licence, http://creativecommons.org/licenses/by/4.0/, as recorded in the arXiv licence metadata for this exact version and shown on the abstract page https://arxiv.org/abs/2311.14953v1. Full licence notice: \texttt{licenses/arxiv-2311.14953-CC-BY-4.0.txt}.\par\smallskip

\noindent\textit{Editorial scope.} Counted unit: the article's lemma l:main (source0.tex lines 77--87). The article's complete original proof of it is delivered with it (source0.tex lines 88--117, one \texttt{proof} environment of 30 lines). No mathematics is written, repaired or re-derived: the statement and the proof are the article's own lines, in the article's own notation, and no retained line is substituted or re-flowed.  No further source-local context is needed: every cross-reference of every retained line resolves inside the two delivered spans.  Nothing was omitted from any retained span, so the package carries no omission record.  No retained line contains a citation, so the wrapper carries no bibliography and no bibliography entry.  No retained line contains a figure environment, an image-inclusion command or drawing code, so no image is delivered and the package declares no asset dependency.  Editorial formatting only, in addition: an AMS \texttt{amsart} preamble that restates the article's own declarations and macros used by the retained lines, namely the environments \texttt{lemma} on separate counters, together with the standard packages those lines need.  The retained environments print in this document as Lemma 1 in order of appearance, whereas the article prints the same items with the numbering of its own full document.  The complete native source of the article as distributed in the arXiv e-print for this exact version is bundled unchanged as \texttt{sources/150456/source0.tex}.
\begin{lemma}\label{l:main}
  Let $f$ and $g$ be non-constant polynomials over $K$. Suppose that
  \begin{itemize}
  \item[(i)] $f$ has a root of multiplicity $a$,
  \item[(ii)] $g$ has a root of multiplicity $b$,
  \item[(iii)] $a$ and $b$ are relatively prime, and
  \item[(iv)] $f(X)-g(y)$ is separable over $K(y)$.
  \end{itemize}
  Then $a$ divides the order of the Galois group of $f(X)-g(y)$ over
  $K(y)$.
\end{lemma}

\begin{proof}
  Let $E$ be a common splitting field of $f$ and $g$ over $K$, and
  $\alpha$ and $\beta$ be the roots of $f$ and $g$ of multiplicity $a$
  and $b$, respectively. As $a$ and $b$ are relatively prime, there
  are positive integers $r$ and $s$ such that $sb-ra=1$. Choose $z$
  with $z^s=y$. Let $E[[z]]$ be the ring of formal power series in $z$
  over $E$ and $E((z))$ be its quotient field, which is the field of
  formal Laurent series in $z$.

  By definition, $K(y)$ is a subfield of $E((z))$, so the Galois group
  of $f(X)-g(y)$ over $E((z))$ is a subgroup of the Galois group of
  $f(X)-g(y)$ over $K(y)$.

  Thus we are done once we know that $f(X)-g(y)$ has an irreducible
  factor over $E((z))$ of degree $a$. Upon replacing $X$ and $y$ with
  $X+\alpha$ and $y+\beta$ we may and do assume that $\alpha=\beta=0$.

  Thus $f(X)=X^au(X)$ and $g(y)=y^bv(y)=z^{sb}v(z^s)$ with
  $u(0)\ne0\ne v(0)$. The degrees of the irreducible factors of
  $f(X)-g(y)$ over $E((z))$ don't change if we replace $X$ with
  $Xz^r$. Recall that $sb-ra=1$, so we are concerned with the
  factorization of\[ H(X)=X^au(Xz^r)-zv(z^s).
  \]
  As $X^a$ and $u(Xz^r)$ are relatively prime, Hensel's Lemma provides
  a factorization $H(X)=A(X)U(X)$ over $E[[z]]$ such that
  $A(X)\equiv X^a\pmod{z}$. % and $U(X)\equiv u(Xz^r)\pmod{z}$.
  Furthermore, $A(0)U(0)=H(0)=-zv(z^s)$, so $z^2$ does not divide
  $A(0)$. Therefore $A(X)$ is an Eisenstein polynomial of degree $a$
  with respect to $z$, and we are done.
\end{proof}

\end{document}
